% \iffalse meta-comment
%
% hit-report-bib.dtx 是开题报告的参考文献页
%
% \fi
%
% \section{开题报告的参考文献页}
%
% \changes{v3.2a}{2026/09/14}{引用样式与著录列表参数抽到 \file{hit-bibstyle.dtx} 共用；
%   \pkg{natbib} 挪到导言区末尾装并先探 \pkg{biblatex}，与终稿同理；正文照终稿排的
%   那一档参考文献另起一页、带标题、进目录}
%
% ^^A ======================================================================
% ^^A Module report-bib: 参考文献加载与 natbib/gbt7714 设置（hit-report）
% ^^A ======================================================================
%
% 底下走 \pkg{natbib}，不带 \cs{bibstyle@inline}，自己定义一份。加载完之后
% 引用样式与著录列表的参数在共享层 \file{hit-bibstyle.dtx}，排在本模块之前。
%
% \pkg{natbib} 挪到导言区末尾装，道理与终稿挪 \pkg{gbt7714} 相同，见
% \file{hit-final-bib.dtx}。\cs{bibstyle@inline} 是普通定义，\cs{bibpunct}
% 到排版时才展开，留在类加载期；\cs{citestyle} 与 \env{thebibliography} 的
% 重定义要等 \pkg{natbib} 装完才轮得到，挂 \texttt{package/natbib/after}。
%
%    \begin{macrocode}
%<*report-bib>
\bool_if:NF \g__hit_final_bool
  {
\AddToHook { begindocument / before }
  {
    \@ifpackageloaded { biblatex }
      { }
      { \RequirePackage[sort&compress,numbers]{natbib} }
  }
\cs_new_protected:Npn \bibstyle@inline
  {
    \bibpunct { \hit@bib@cite@open: } { \hit@bib@cite@close: }
      { \hit@bib@cite@sep: } {n} {,} {，}
  }
%    \end{macrocode}
%
% \changes{v3.1c}{2023/04/16}{深圳校区硕士开题部分参考文献没有序号。\issue[141]}
% 深圳校区硕士开题部分参考文献没有序号。其他校区的提了问题再改。
%
% \emph{正文照终稿排的那一档，排在后置部分的参考文献自带标题。} 深圳本科那
% 两份报告的参考文献在 iota-hit 对原件的复刻件里是一整页：标题落在版心顶起
% 30.23（与正文每章的标题同数），条目步进 19.45，目录里也有这一条。这边先前
% 只给深圳硕士开题发，深圳本科的参考文献就直接接在上一节末尾往下排。
% 另起一页不用另发，顶级标题前面已经无条件换页，见
% \file{src/flow/hit-report-matter.dtx}。
%
% \emph{判据要带上「排在后置部分」。} 表单自己列了参考文献那一节的（示例里
% 深圳博士那份的「九、主要参考文献」）是在正文里写 \cs{bibliography}，标题
% 用户自己写了；再自动发一个就成了两页标题。所以只认
% \cs{hitbackmatter} 里那一次，也就是 \option{struct} 的 \option{bibliography}
% 部件。写在正文里的照旧不发。
%
% 标题按语种取：中文版「参考文献」，英文版「References」。先前写死中文那一条。
% \begin{environment}{thebibliography}
%    \begin{macrocode}
\cs_new:Npn \hit@report@bibliography@name
  {
    \bool_if:NTF \g__hit_en_bool
      { \hit@term@bibliography@heading@en }
      { \hit@term@bibliography@heading@zh }
  }
\cs_new_protected:Npn \hit@report@bibliography@heading
  {
    \hit@report@top@heading * { \hit@report@bibliography@name }
    \addcontentsline { toc } { section } { \hit@report@bibliography@name }
  }
\AddToHook { package / natbib / after }
  {
\citestyle{numerical}
\RenewDocumentEnvironment{thebibliography}{m}{%
  \bool_lazy_or:nnT
    {
      \bool_lazy_and_p:nn
        { \bool_if_p:N \g__hit_final_body_bool }
        { \bool_if_p:N \g__hit_matter_back_bool }
    }
    {
      \bool_lazy_all_p:n
        {
          { \bool_if_p:N \g__hit_shenzhen_bool }
          { \bool_if_p:N \g__hit_master_bool }
          { \bool_if_p:N \g__hit_proposal_bool }
        }
    }
    { \hit@report@bibliography@heading }
  \hit@bib@list:nn {#1} { \hit@bib@labelsep: }
  \sloppy\frenchspacing
  \flushbottom
  \clubpenalty0
  \@clubpenalty \clubpenalty
  \widowpenalty0%
  \interlinepenalty-50%
  \sfcode`\.\@m
}
{\hit@if@shenzhen@doctor@interim@TF
  { }
  {\cs_set_nopar:Npn \@noitemerr
  {\@latex@warning{Empty `thebibliography' environment}}%
\endlist\frenchspacing}}
%    \end{macrocode}
%
%    \begin{macrocode}
\providecommand \printbibliography { \hit@bib@print: }
\providecommand \supercite { \cite }
  }
%    \end{macrocode}
%
% \begin{macro}{\bibauthor}
% \cs{bibauthor}\marg{作者表} 由 \file{.bst} 调用，按 \option{bibmaxauthor} 截断，
% 按 \option{bibauthorupper} 决定要不要全大写。两个类各留一份而没有并进
% \file{hit-bibstyle.dtx}：单个作者的排法用 \cs{exp\_args:NNe} 在定义时就把
% \option{bibauthorupper} 的取值展开定死，得排在 \cs{ProcessKeyOptions} 之后，
% 而共享层排在类骨架之前。
%    \begin{macrocode}
\exp_args:NNe \cs_new:Nn \g_bib_author_cs:n {
    \bool_if:NTF \g_bib_authorupper_bool {
        \exp_not:N \text_uppercase:n {#1}
    } {
        {#1}
    }
}
\NewDocumentCommand{\bibauthor}{m}{
  \int_zero:N \l_tmpa_int
  \exp_args:Ne \tl_map_inline:nn {
    \exp_args:Ne \tl_tail:n {#1}
  } {
    \int_compare:nNnT
      \l_tmpa_int = \g_bib_maxauthor_int {
      \prg_break:n {\tl_head:N {#1}}
    }
    \int_incr:N \l_tmpa_int
    \g_bib_author_cs:n {##1}
  }
  \prg_break_point:
}
%    \end{macrocode}
% \end{macro}
%
% \subsection{biblatex 一线}
%
% 与终稿同一套机制（见 \file{hit-final-bib.dtx} 与
% \file{hit-bibstyle.dtx}），排版参数照报告自己的 \env{thebibliography} 搬：
% 标号后间距 0.5em，分页惩罚归零，深圳硕士开题才排那行节标题。报告
% 不用全角序号，\texttt{gbbiblabel} 保持默认的半角方括号。
%
% \begin{macro}{\hit@bib@biber@heading:,\hit@bib@biber@env@begin:,\hit@bib@biblatex@setup:}
%    \begin{macrocode}
\cs_new_protected:Npn \hit@bib@biber@heading:
  {
    \bool_lazy_all:nT
      {
        { \bool_if_p:N \g__hit_shenzhen_bool }
        { \bool_if_p:N \g__hit_master_bool }
        { \bool_if_p:N \g__hit_proposal_bool }
      }
      { \hit@report@bibliography@heading }
  }
\cs_new_protected:Npn \hit@bib@biber@env@begin:
  {
    \list
      {
        \printtext [ labelnumberwidth ]
          { \printfield { labelprefix } \printfield { labelnumber } }
      }
      {
        \str_if_eq:VnT \g_bib_labelalign_str { left }
          { \RenewDocumentCommand \makelabel { m } { ##1 \hfill } }
        \dim_set:Nn \labelwidth { \labelnumberwidth }
        \dim_set:Nn \labelsep { \hit@bib@labelsep: }
        \dim_set:Nn \itemindent { 0pt }
        \dim_set:Nn \leftmargin { \labelsep + \labelwidth }
        \skip_zero:N \itemsep
        \skip_zero:N \parsep
      }
    \sloppy \frenchspacing
    \flushbottom
    \clubpenalty 0
    \@clubpenalty \clubpenalty
    \widowpenalty 0
    \interlinepenalty -50
    \sfcode `\. \@m
  }
\cs_new:cpn { gbbiblabelopt@hithalfwidth }
  { \renewcommand { \mkgbnumlabel } [1] { \mbox { [##1] } } }
\cs_new_protected:Npn \hit@bib@biblatex@setup:
  {
    \cs_if_exist:NT \mkgbnumlabel
      { \ExecuteBibliographyOptions { gbbiblabel = hithalfwidth } }
    \defbibheading { bibliography } [ \bibname ] { \hit@bib@biber@heading: }
    \defbibenvironment { bibliography }
      { \hit@bib@biber@env@begin: } { \endlist \frenchspacing } { \item }
  }
%    \end{macrocode}
% \end{macro}
%
%    \begin{macrocode}
  }
%</report-bib>
%    \end{macrocode}
% \end{environment}
%
% \Finale
